%% =============================================================================
%% chapters/02-business-layer.tex
%% ArchiMate Business & Clinical Process Architecture Layer
%% =============================================================================

\chapter{Business \& Clinical Architecture}
\label{chap:business_layer}

The Business Layer specifies the healthcare business environment, defining external business actors, internal operational roles, business services exposed to healthcare facilities, business processes, and the core business objects exchanged across clinical pathways.

\begin{figure}[htbp]
    \centering
    \begin{adjustbox}{max width=\textwidth}
        %% =============================================================================
%% diagrams/fig-clinical-process.tex
%% ArchiMate Business & Clinical Architecture View
%% =============================================================================
\begin{tikzpicture}[node distance=1.0cm and 0.8cm, auto]

    % --- Background Layer Plates ---
    \draw[archi-plate-bg] (-0.8, 6.2) rectangle (15.8, 9.8);
    \node[archi-plate-title] at (-0.6, 9.6) {Business Actors \& Roles};

    \draw[archi-plate-bg] (-0.8, 4.2) rectangle (15.8, 6.0);
    \node[archi-plate-title] at (-0.6, 5.8) {Clinical Business Services};

    \draw[archi-plate-bg] (-0.8, 2.2) rectangle (15.8, 4.0);
    \node[archi-plate-title] at (-0.6, 3.8) {End-to-End Business Process Sequence};

    \draw[archi-plate-bg] (-0.8, 0.2) rectangle (15.8, 2.0);
    \node[archi-plate-title] at (-0.6, 1.8) {Clinical Business Objects \& Information};

    % --- Business Actors & Roles ---
    \node[archi-business-actor, minimum width=3.4cm] (act_hospital)  at (1.0, 8.5) {\archibadge{Business Actor}\archiname{Hospital EHR}\\\archidesc{Clinical Source System}};
    \node[archi-business-actor, minimum width=3.4cm] (act_lab)       at (5.5, 8.5) {\archibadge{Business Actor}\archiname{Diagnostic Lab}\\\archidesc{LIS / Pathology System}};
    \node[archi-business-actor, minimum width=3.4cm] (act_clinician) at (12.5, 8.5) {\archibadge{Business Actor}\archiname{Practitioner}\\\archidesc{Clinician / Care Team}};

    \node[archi-business-role, minimum width=4.5cm]  (rol_producer)  at (3.25, 7.0) {\archibadge{Business Role}\archiname{Clinical Data Producer}\\\archidesc{HL7 Sender / Ingestion Agent}};
    \node[archi-business-role, minimum width=4.0cm]  (rol_consumer)  at (12.5, 7.0) {\archibadge{Business Role}\archiname{Clinical Care Provider}\\\archidesc{FHIR Consumer / Analyst}};

    % --- Business Services ---
    \node[archi-business-service, minimum width=3.8cm] (srv_adt)   at (1.2, 5.0) {\archibadge{Business Service}\archiname{Admission Service}\\\archidesc{HL7 ADT Gateway}};
    \node[archi-business-service, minimum width=3.8cm] (srv_mfn)   at (5.6, 5.0) {\archibadge{Business Service}\archiname{Directory Service}\\\archidesc{HL7 MFN Ingestion}};
    \node[archi-business-service, minimum width=4.2cm] (srv_query) at (12.5, 5.0) {\archibadge{Business Service}\archiname{FHIR Clinical Access}\\\archidesc{Resource Discovery \& REST}};

    % --- Business Processes ---
    \node[archi-business-process, minimum width=3.5cm] (prc_ingest)  at (0.9, 3.0) {\archibadge{Business Process}\archiname{1. Ingest Trigger}\\\archidesc{Capture \& Verify Payload}};
    \node[archi-business-process, minimum width=3.5cm] (prc_match)   at (5.1, 3.0) {\archibadge{Business Process}\archiname{2. Reconcile}\\\archidesc{EMPI Resolution \& ID}};
    \node[archi-business-process, minimum width=3.5cm] (prc_persist) at (9.3, 3.0) {\archibadge{Business Process}\archiname{3. Store Record}\\\archidesc{Canonical Persistence}};
    \node[archi-business-process, minimum width=3.5cm] (prc_serve)   at (13.5, 3.0) {\archibadge{Business Process}\archiname{4. Expose Resource}\\\archidesc{FHIR Query \& Subscriptions}};

    % --- Business Objects ---
    \node[archi-business-object, minimum width=3.5cm] (obj_hl7)     at (0.9, 1.0) {\archibadge{Business Object}\archiname{HL7 v2 Message}\\\archidesc{ADT / MFN Frame}};
    \node[archi-business-object, minimum width=5.5cm] (obj_patient) at (6.8, 1.0) {\archibadge{Business Object}\archiname{Clinical Entity Model}\\\archidesc{FHIR Patient, Practitioner, Provenance}};
    \node[archi-business-object, minimum width=3.5cm] (obj_audit)   at (13.5, 1.0) {\archibadge{Business Object}\archiname{AuditEvent Record}\\\archidesc{Audit Trail \& Safety}};

    % --- Relationships ---
    % Actor -> Role (Assignment)
    \draw[archi-assignment] (act_hospital) -- (rol_producer);
    \draw[archi-assignment] (act_lab) -- (rol_producer);
    \draw[archi-assignment] (act_clinician) -- (rol_consumer);

    % Role -> Service (Serving)
    \draw[archi-serving] (srv_adt) -- (rol_producer);
    \draw[archi-serving] (srv_mfn) -- (rol_producer);
    \draw[archi-serving] (srv_query) -- (rol_consumer);

    % Process Realizes Service
    \draw[archi-realization] (prc_ingest) -- (srv_adt);
    \draw[archi-realization] (prc_match) -- (srv_adt);
    \draw[archi-realization] (prc_persist) -- (srv_adt);
    \draw[archi-realization] (prc_serve) -- (srv_query);

    % Process Flow (Triggering)
    \draw[archi-triggering] (prc_ingest) -- (prc_match);
    \draw[archi-triggering] (prc_match) -- (prc_persist);
    \draw[archi-triggering] (prc_persist) -- (prc_serve);

    % Process Accesses Business Objects
    \draw[archi-access] (prc_ingest) -- (obj_hl7);
    \draw[archi-access] (prc_match) -- (obj_patient);
    \draw[archi-access] (prc_persist) -- (obj_patient);
    \draw[archi-access] (prc_persist) -- (obj_audit);
    \draw[archi-access] (prc_serve) -- (obj_patient);

\end{tikzpicture}

    \end{adjustbox}
    \caption{ArchiMate Business and Clinical Architecture View for Harmonia HIE}
    \label{fig:clinical_process}
\end{figure}

\section{Healthcare Business Actors \& Operational Roles}

Harmonia interacts with a diverse ecosystem of clinical actors, abstracting their interactions through formal business roles:

\begin{table}[htbp]
\centering
\small
\begin{tabularx}{\textwidth}{l p{4.0cm} X}
\toprule
\textbf{Business Actor} & \textbf{Assigned Business Role} & \textbf{Healthcare Operational Scope} \\
\midrule
\textbf{Hospital EHR / PAS} & \textit{Clinical Data Producer} & Sends continuous real-time HL7 v2 ADT feeds (admit, discharge, transfer, demographic updates) over MLLP connections. \\
\addlinespace
\textbf{Diagnostic Laboratory} & \textit{Diagnostic Result Producer} & Dispatches lab order notifications and observation reports (HL7 ORU / FHIR \texttt{Observation}) to the platform. \\
\addlinespace
\textbf{External System (HIS / LIS)} & \textit{Clinical Data Consumer} & Receives outbound HL7 v2 triggers, ADT notifications, and MFN directory updates via MLLP connections and returns synchronous acknowledgments (\texttt{ACK}). \\
\addlinespace
\textbf{Practitioner / Clinician} & \textit{Clinical Care Provider} & Queries longitudinal patient records, reviews historical clinical encounters, inspects diagnostic summaries, and looks up provider directory entries via web UI or FHIR APIs. \\
\addlinespace
\textbf{Directory Steward} & \textit{Directory Governance Authority} & Governs healthcare provider directory entries, verifies practitioner credentialing, oversees organizational hierarchies, and reviews flagged directory change requests. \\
\addlinespace
\textbf{Credentialing Authority} & \textit{Authoritative Identity Source} & National or state medical boards (e.g. AHPRA, HI Service) publishing authoritative provider identifiers (HPI-I, HPI-O) and qualification updates. \\
\addlinespace
\textbf{Health Network Admin} & \textit{Operations Administrator} & Monitors message queue depths, inspects Artemis cluster topologies, triggers task sequence reload, and audits failed pipelines. \\
\bottomrule
\end{tabularx}
\caption{Harmonia Business Actors, Assigned Roles, and Operational Responsibilities}
\label{tab:business_actors}
\end{table}

\section{Clinical Business Services}

Harmonia exposes six core business services to external healthcare organizations:

\subsection{1. Patient Admission \& Movement Service (\texttt{HL7-ADT})}
Accepts real-time patient admission, discharge, and transfer (ADT) event streams adhering to HL7 v2.3--v2.5 standards across standard trigger events:
\begin{itemize}
    \item \textbf{ADT\^{}A01 / A04}: Patient Admission / Outpatient Registration.
    \item \textbf{ADT\^{}A03}: Patient Discharge / End of Encounter.
    \item \textbf{ADT\^{}A08}: Patient Demographics / Record Update.
    \item \textbf{ADT\^{}A40}: Merge Patient Identifiers (MRN / EMPI consolidation).
\end{itemize}

\subsection{2. Practitioner \& Staff Directory Service (\texttt{HL7-MFN})}
Maintains an authoritative, synchronized regional directory of healthcare practitioners, specialties, role assignments, and organizational facility hierarchies via HL7 MFN master file notifications.

\subsection{3. FHIR R5 Clinical Data Discovery Service}
Provides RESTful search, retrieval, and subscription capabilities across 15 FHIR R5 resource types, enabling authorized clinical applications to query patient histories, encounters, and audit logs.

\subsection{4. Operations Telemetry \& Workflow Administration Service}
Exposes real-time health metrics, ActiveMQ Artemis broker statistics, Infinispan cache hit rates, and Ponos task execution statuses to health network operators.

\subsection{5. Clinical Outbound Notification \& Egress Service (\texttt{HL7-MLLP-OUT})}
Dispatches transformed, enriched clinical trigger events (ADT, MFN, ORU) and outbound clinical notifications to remote hospital information systems and diagnostic endpoints over resilient MLLP/TCP connections, enforcing synchronous HL7 ACK/NACK validation, automatic retries, and audit tracking.

\subsection{6. FHIR R5 Provider Directory Management Service}
Exposes synchronous read/search and governed asynchronous change management APIs for healthcare directory entities (\texttt{Practitioner}, \texttt{PractitionerRole}, \texttt{Organization}, \texttt{Location}, \texttt{HealthcareService}, \texttt{Endpoint}, \texttt{Group}). It enforces automated referential verification, optimistic concurrency locking, duplicate identifier protection, and auditable FHIR \texttt{Task} tracking.

\section{Clinical Business Process Workflows}

\subsection{Process: Ingest and Reconcile HL7 ADT Trigger}
The primary clinical ingestion workflow executes through five discrete business stages:
\begin{enumerate}
    \item \textbf{Stage 1: Capture MLLP Frame}: Inbound MLLP gateway receives the TCP stream, validates the MLLP envelope framing (\texttt{<VT>} and \texttt{<FS><CR>}), and returns an immediate HL7 \texttt{ACK} to release the source hospital system.
    \item \textbf{Stage 2: Envelope Encapsulation}: The raw message is wrapped into a \texttt{PetasosMessage} and published onto the clustered Petasos queue (\texttt{clinical.tasks.inbound}).
    \item \textbf{Stage 3: Patient Identity Resolution (EMPI)}: The \texttt{PatientIdentityUpdateErgon} activity extracts the Medical Record Number (MRN), queries the Enterprise Master Patient Index (EMPI) in Mneme cache, and assigns a canonical patient identifier.
    \item \textbf{Stage 4: Demographics Merge \& Enrichment}: The \texttt{PatientDemographicsUpdateErgon} merges telecom, residential addresses, contact persons, and marital status into a consolidated FHIR \texttt{Patient} resource.
    \item \textbf{Stage 5: Audit \& Persistence}: A \texttt{Provenance} and \texttt{AuditEvent} record is linked to the transaction, and the resulting resources are stored in Mneme cache and Mnemosyne PostgreSQL.
\end{enumerate}

\subsection{Process: Outbound MLLP Dispatch and ACK Validation Lifecycle}
The outbound clinical egress workflow executes through five structured business stages:
\begin{enumerate}
    \item \textbf{Stage 1: Egress Queue Enqueue}: The Ponos WorkEngine or upstream routing activity publishes an outbound dispatch task to the endpoint's dedicated Petasos queue (\texttt{petasos.queue.mllp.outbound.<endpoint-id>}).
    \item \textbf{Stage 2: Endpoint Resolution \& Dequeue}: The dedicated \texttt{pylai-mllp-out} instance consumes the task and verifies target destination parameters (host, port, timeouts) via the dynamic \texttt{MllpDestinationRegistry}.
    \item \textbf{Stage 3: MLLP Framing \& Wire Transmission}: The Apache Camel MLLP client component frames the HL7 v2 payload with standard delimiters (\texttt{<VT>} and \texttt{<FS><CR>}) and transmits the message over a dedicated TCP socket.
    \item \textbf{Stage 4: Synchronous ACK / NACK Evaluation}: The gateway awaits the remote system's acknowledgment frame, parsing the \texttt{MSH} and \texttt{MSA} segments to validate acknowledgment codes (\texttt{AA}/\texttt{CA} for success vs \texttt{AE}/\texttt{AR}/\texttt{CE}/\texttt{CR} for application/commit errors).
    \item \textbf{Stage 5: State Machine Finalization \& Audit}: The task status is updated in the Mneme cache grid, recording FHIR \texttt{Communication}, \texttt{Task} (\texttt{Pragma}), and \texttt{Provenance} records with remote acknowledgment details and transmission metrics.
\end{enumerate}

\subsection{Process: Governed Provider Registry Change Ingestion \& Verification}
The governed provider directory change workflow executes through five structured business stages:
\begin{enumerate}
    \item \textbf{Stage 1: Ingestion \& Pragma Generation}: The \texttt{pylai-fhir-registry} gateway receives a \texttt{POST} or \texttt{PUT} change request, performs structural validation, generates a \texttt{ProviderRegistryChangePragma}, and dispatches a \texttt{PetasosMessage} to Artemis queue \texttt{harmonia.provider.registry.change.request}, returning \texttt{HTTP 202 Accepted} with \texttt{Location: /Task/\{id\}}.
    \item \textbf{Stage 2: Sequence Orchestration}: The Ponos WorkEngine consumes the change request and triggers the Praxis sequence \texttt{seq-provider-registry-change-pipeline}, routing the task to the dedicated resource Ergon activity.
    \item \textbf{Stage 3: Referential \& Business Rule Verification}: The Ergon activity queries \texttt{ProviderRegistryReferenceValidator} to verify that all referenced target entities (Practitioner, Organization, Location, HealthcareService, Endpoint) exist and are active in Mnemosyne storage, while verifying identifier uniqueness.
    \item \textbf{Stage 4: Automated Approval \& Versioned Persistence}: If verification passes, the change is approved and committed to Mnemosyne PostgreSQL (\texttt{hie\_fhir\_resources}), incrementing the resource \texttt{version\_id} and updating \texttt{last\_updated}. If verification fails, the request transitions to \texttt{REJECTED} with an embedded FHIR \texttt{OperationOutcome}.
    \item \textbf{Stage 5: Task Status Reflection}: The underlying \texttt{Pragma} transitions to \texttt{COMPLETED} (or \texttt{REJECTED}), making the updated resource immediately visible on synchronous \texttt{GET} read/search paths and reflecting completion on \texttt{GET /Task/\{id\}}.
\end{enumerate}

\section{Core Business Objects}

\begin{description}
    \item[\textbf{HL7 v2 Message Stream}] The raw wire-format clinical transmission received from legacy hospital interface engines.
    \item[\textbf{Canonical Patient Record (\texttt{FHIR::Patient})}] The consolidated, deduplicated clinical identity representing the care recipient.
    \item[\textbf{Task State Record (\texttt{Pragma})}] The authoritative audit representation of an in-flight integration task, including all historical state checkpoints.
    \item[\textbf{Provider Registry Change Pragma}] Specialized canonical task entity capturing directory change requests, target resources, correlation IDs, and validation outcomes.
    \item[\textbf{Authoritative Provider Graph}] Interconnected FHIR R5 directory resources representing practitioners, roles, organizations, locations, healthcare services, endpoints, and groups.
    \item[\textbf{Audit \& Provenance Trail (\texttt{FHIR::AuditEvent} / \texttt{FHIR::Provenance})}] Immutable compliance records establishing source attribution, timestamping, and processing integrity.
\end{description}

\newpage
